\documentclass[11pt,a4paper]{article}

\usepackage[margin=1in]{geometry}
\usepackage[T1]{fontenc}
\usepackage[utf8]{inputenc}
\usepackage{lmodern}
\usepackage{xcolor}
\usepackage{booktabs}
\usepackage{longtable}
\usepackage{array}
\usepackage{enumitem}
\usepackage{amsmath}
\usepackage{amssymb}
\usepackage{listings}
\usepackage[hidelinks,breaklinks=true]{hyperref}

\definecolor{codebg}{RGB}{248,248,248}
\definecolor{codeframe}{RGB}{220,220,220}
\definecolor{commentgreen}{RGB}{0,110,55}
\definecolor{keywordblue}{RGB}{0,70,150}
\definecolor{stringred}{RGB}{150,35,35}

\lstdefinestyle{cppstyle}{
  language=C++,
  basicstyle=\ttfamily\scriptsize,
  keywordstyle=\color{keywordblue}\bfseries,
  commentstyle=\color{commentgreen},
  stringstyle=\color{stringred},
  numbers=left,
  numberstyle=\tiny\color{gray},
  stepnumber=1,
  numbersep=6pt,
  breaklines=true,
  breakatwhitespace=false,
  tabsize=2,
  keepspaces=true,
  showstringspaces=false,
  backgroundcolor=\color{codebg},
  frame=single,
  rulecolor=\color{codeframe},
  columns=fullflexible,
  captionpos=b
}

\lstdefinestyle{cmdstyle}{
  basicstyle=\ttfamily\scriptsize,
  numbers=left,
  numberstyle=\tiny\color{gray},
  numbersep=6pt,
  breaklines=true,
  keepspaces=true,
  showstringspaces=false,
  backgroundcolor=\color{codebg},
  frame=single,
  rulecolor=\color{codeframe},
  columns=fullflexible,
  captionpos=b
}

\newcommand{\codepath}[1]{\path{#1}}

\title{VANET Security Implementation Report}
\author{VANET ns-3 Security Research Workspace}
\date{April 29, 2026}

\begin{document}
\maketitle

\begin{abstract}
This report summarizes the implemented VANET security flow in the ns-3
\codepath{vanet-security} module. The focus is the five main security stages:
system setup, pseudonym registration, KGC partial-key generation, vehicle
signature generation, and relay-side aggregate verification. Each stage is
documented with formulas and the exact source-code excerpts that implement it.
\end{abstract}

\tableofcontents
\newpage

\section{Implementation Scope}

The implementation lives mainly in:

\begin{longtable}{p{0.35\linewidth}p{0.55\linewidth}}
\toprule
\textbf{File} & \textbf{Role} \\
\midrule
\codepath{src/vanet-security/model/pbc-crypto.cc} & PBC setup, PID generation, partial key, signing, aggregation, verification. \\
\codepath{src/vanet-security/model/ecc-crypto.cc} & ECDSA P-256 key generation, signing, verification, SHA-256 helper. \\
\codepath{src/vanet-security/model/vehicle-app.cc} & Vehicle registration, credential storage, V2V signing. \\
\codepath{src/vanet-security/model/ta-app.cc} & TA certificate signing and PID2 generation. \\
\codepath{src/vanet-security/model/kgc-app.cc} & KGC key response and partial private key issue. \\
\codepath{src/vanet-security/model/bs-app.cc} & BS orchestration between vehicle, RSU, TA, and KGC. \\
\codepath{src/vanet-security/model/bs-relay-app.cc} & LTE/NR relay batching and aggregate verification. \\
\codepath{scratch/vanet-security-lte.cc} & LTE scenario integration and \codepath{--securityMode}. \\
\codepath{scratch/vanet-security-nr.cc} & NR/5G-LENA scenario integration and \codepath{--securityMode}. \\
\bottomrule
\end{longtable}

\section{Security Modes}

The LTE and NR scenarios support:

\begin{itemize}[leftmargin=*]
  \item \textbf{PBC mode}: \codepath{--securityMode=pbc}. Enables PBC on TA, KGC, vehicles, and BS relay. The relay verifies V2V packets using aggregate pairing verification.
  \item \textbf{ECC mode}: \codepath{--securityMode=ecc}. Disables PBC at runtime and keeps the ECDSA certificate/signature baseline.
\end{itemize}

\section{High-Level Message Flow}

\begin{longtable}{p{0.20\linewidth}p{0.17\linewidth}p{0.17\linewidth}p{0.36\linewidth}}
\toprule
\textbf{Message} & \textbf{From} & \textbf{To} & \textbf{Purpose} \\
\midrule
\codepath{REGISTER_REQ} & Vehicle & RSU/BS & Vehicle requests credentials and sends PID1 in PBC mode. \\
\codepath{KEY_REQ} & BS & KGC & BS asks KGC for ECDSA keys and optional PBC partial key. \\
\codepath{KEY_RESP} & KGC & BS & KGC returns key material and optional \(S_i\). \\
\codepath{CERT_REQ} & BS & TA & BS asks TA for certificate and PID2. \\
\codepath{CERT_RESP} & TA & BS & TA returns certificate, PID2, and timestamp \(T_i\). \\
\codepath{REGISTER_RESP} & BS/RSU & Vehicle & Vehicle receives ECC and PBC credentials. \\
\codepath{V2V_PBC_DATA} & Vehicle & BS relay & PBC-authenticated V2V message for relay verification. \\
\bottomrule
\end{longtable}

\section{Notation}

\begin{longtable}{p{0.18\linewidth}p{0.72\linewidth}}
\toprule
\textbf{Symbol} & \textbf{Meaning} \\
\midrule
\(G_1\) & Additive elliptic-curve group used by PBC. \\
\(G_T\) & Multiplicative target group of the pairing. \\
\(P\) & Generator of \(G_1\). \\
\(e(\cdot,\cdot)\) & Bilinear pairing \(G_1\times G_1\rightarrow G_T\). \\
\(s\) & KGC master secret. \\
\(P_{pub}=sP\) & KGC public key. \\
\(t\) & TA master secret. \\
\(T_{pub}=tP\) & TA public key. \\
\(\mathrm{PID}_i\) & Vehicle pseudonym tuple \((\mathrm{PID1}_i,\mathrm{PID2}_i,T_i)\). \\
\(S_i\) & KGC-issued partial private key for vehicle \(i\). \\
\(x_i,Q_i\) & Vehicle secret/public verification key, \(Q_i=x_iP\). \\
\(w_i,W_i\) & Per-message random scalar/point, \(W_i=w_iP\). \\
\(\theta\) & Shared state string, default \codepath{traffic-alert}. \\
\(\psi_i\) & PBC signature component. \\
\bottomrule
\end{longtable}

\section{Stage 1: System Setup}

\subsection{Goal}

The TA and KGC initialize the pairing system and publish the public parameters
needed by registration, signing, and verification.

\subsection{Formulas}

\[
s\xleftarrow{\$}\mathbb{Z}_r,\qquad P_{pub}=sP
\]

\[
t\xleftarrow{\$}\mathbb{Z}_r,\qquad T_{pub}=tP
\]

\[
\theta=\text{traffic-alert}
\]

\subsection{Implementation}

\begin{itemize}[leftmargin=*]
  \item \codepath{PbcCrypto::SetupKgc()} creates \(P\), samples \(s\), and computes \(P_{pub}=sP\).
  \item \codepath{PbcCrypto::SetupTa()} samples \(t\) and computes \(T_{pub}=tP\).
  \item LTE/NR scenario code enables or disables PBC using \codepath{--securityMode=pbc|ecc}.
  \item In PBC mode, the scenario injects KGC \(P_{pub}\) into the BS relay so the relay can verify aggregate signatures.
\end{itemize}

\lstinputlisting[
  style=cppstyle,
  firstline=70,
  lastline=143,
  caption={Stage 1 code: PBC setup for KGC and TA.},
  label={lst:stage1-pbc-setup}
]{simulation_workspace/ns-allinone-3.44/ns-3.44/src/vanet-security/model/pbc-crypto.cc}

\lstinputlisting[
  style=cppstyle,
  firstline=211,
  lastline=252,
  caption={Stage 1 code: LTE runtime security-mode switch.},
  label={lst:stage1-lte-mode}
]{simulation_workspace/ns-allinone-3.44/ns-3.44/scratch/vanet-security-lte.cc}

\lstinputlisting[
  style=cppstyle,
  firstline=442,
  lastline=504,
  caption={Stage 1 code: LTE installs PBC-enabled apps and injects Ppub.},
  label={lst:stage1-lte-enable}
]{simulation_workspace/ns-allinone-3.44/ns-3.44/scratch/vanet-security-lte.cc}

\section{Stage 2: Vehicle Registration and Pseudonym Generation}

\subsection{Goal}

The vehicle obtains an active pseudonym. This pseudonym is later used in the
PBC signing equation instead of relying directly on the plain node ID.

\subsection{Formulas}

The vehicle creates:

\[
d_i\xleftarrow{\$}\mathbb{Z}_r,\qquad \mathrm{PID1}_i=d_iP
\]

The TA computes:

\[
q_i=t\cdot \mathrm{PID1}_i
\]

\[
h_i=\operatorname{SHA256}(q_i\parallel T_{pub})
\]

\[
\mathrm{PID2}_i=\operatorname{IDBytes}_i\oplus h_i[0..3]
\]

The active pseudonym is:

\[
\mathrm{PID}_i=(\mathrm{PID1}_i,\mathrm{PID2}_i,T_i)
\]

\subsection{Implementation}

\begin{itemize}[leftmargin=*]
  \item \codepath{VehicleApp::SendRegisterRequest()} sends the vehicle ID and PID1.
  \item \codepath{PbcCrypto::GeneratePid1()} implements \(\mathrm{PID1}_i=d_iP\).
  \item \codepath{TaApp::HandleCertReq()} signs the ECC certificate and asks PBC code to compute PID2.
  \item \codepath{PbcCrypto::ComputePid2()} implements the TA-side PID2 formula.
\end{itemize}

\lstinputlisting[
  style=cppstyle,
  firstline=115,
  lastline=217,
  caption={Stage 2 code: vehicle sends registration request with PID1.},
  label={lst:stage2-register}
]{simulation_workspace/ns-allinone-3.44/ns-3.44/src/vanet-security/model/vehicle-app.cc}

\lstinputlisting[
  style=cppstyle,
  firstline=158,
  lastline=260,
  caption={Stage 2 code: PID1 and PID2 generation.},
  label={lst:stage2-pid}
]{simulation_workspace/ns-allinone-3.44/ns-3.44/src/vanet-security/model/pbc-crypto.cc}

\lstinputlisting[
  style=cppstyle,
  firstline=153,
  lastline=246,
  caption={Stage 2 code: TA returns certificate, PID2, and timestamp.},
  label={lst:stage2-ta}
]{simulation_workspace/ns-allinone-3.44/ns-3.44/src/vanet-security/model/ta-app.cc}

\section{Stage 3: KGC Partial Private Key Generation}

\subsection{Goal}

The KGC creates the partial private key \(S_i\), binding the active vehicle
pseudonym to the KGC master secret \(s\).

\subsection{Formulas}

Encode the active pseudonym:

\[
\operatorname{PIDBytes}_i=
\operatorname{len}(\mathrm{PID1}_i)\parallel \mathrm{PID1}_i
\parallel \operatorname{len}(\mathrm{PID2}_i)\parallel \mathrm{PID2}_i
\parallel T_i
\]

Hash the pseudonym to \(G_1\):

\[
\Pi_i=H_1(\operatorname{PIDBytes}_i\parallel P_{pub})
\]

Compute the partial key:

\[
S_i=s\Pi_i
\]

Freshness check:

\[
0\le T_r-T_i\le \Delta_{\text{fresh}}
\]

\subsection{Implementation}

\begin{itemize}[leftmargin=*]
  \item \codepath{BsApp} sends PID1, PID2, and timestamp to KGC after TA data is available.
  \item \codepath{KgcApp::HandleKeyReq()} checks timestamp freshness before issuing \(S_i\).
  \item \codepath{PbcCrypto::EncodePid()} builds the exact byte representation of \(\mathrm{PID}_i\).
  \item \codepath{PbcCrypto::ComputePartialKey()} computes \(S_i=s\Pi_i\).
\end{itemize}

\lstinputlisting[
  style=cppstyle,
  firstline=121,
  lastline=228,
  caption={Stage 3 code: KGC handles key request and partial-key issue.},
  label={lst:stage3-kgc}
]{simulation_workspace/ns-allinone-3.44/ns-3.44/src/vanet-security/model/kgc-app.cc}

\lstinputlisting[
  style=cppstyle,
  firstline=276,
  lastline=361,
  caption={Stage 3 code: PID encoding and partial-key calculation.},
  label={lst:stage3-partial}
]{simulation_workspace/ns-allinone-3.44/ns-3.44/src/vanet-security/model/pbc-crypto.cc}

\lstinputlisting[
  style=cppstyle,
  firstline=164,
  lastline=303,
  caption={Stage 3 code: BS requests and stores KGC response material.},
  label={lst:stage3-bs}
]{simulation_workspace/ns-allinone-3.44/ns-3.44/src/vanet-security/model/bs-app.cc}

\section{Stage 4: Vehicle Full Key and Signature Generation}

\subsection{Goal}

The vehicle uses \(S_i\), its own secret \(x_i\), and fresh randomness \(w_i\)
to sign a V2V alert message.

\subsection{Formulas}

Vehicle full-key component:

\[
x_i\xleftarrow{\$}\mathbb{Z}_r,\qquad Q_i=x_iP
\]

Per-message random point:

\[
w_i\xleftarrow{\$}\mathbb{Z}_r,\qquad W_i=w_iP
\]

State and message hashes:

\[
X=H_1(\theta)
\]

\[
R_i=H_1(\theta\parallel M_i\parallel \operatorname{PIDBytes}_i\parallel Q_i\parallel W_i)
\]

Signature:

\[
\psi_i=S_i+x_iX+w_iR_i
\]

Transmitted PBC authentication tuple:

\[
\sigma_i=(\mathrm{PID1}_i,\mathrm{PID2}_i,T_i,Q_i,W_i,\psi_i,\theta)
\]

\subsection{Implementation}

\begin{itemize}[leftmargin=*]
  \item \codepath{VehicleApp::HandleRsuRead()} stores PID2, \(T_i\), \(S_i\), and generates \(Q_i=x_iP\).
  \item \codepath{PbcCrypto::SignMessage()} computes \(W_i\), \(X\), \(R_i\), and \(\psi_i\).
  \item \codepath{VehicleApp::SendV2vMessage()} selects PBC signing only when PBC and BS relay are enabled.
\end{itemize}

\lstinputlisting[
  style=cppstyle,
  firstline=248,
  lastline=356,
  caption={Stage 4 code: vehicle stores PBC credentials and generates Qi.},
  label={lst:stage4-creds}
]{simulation_workspace/ns-allinone-3.44/ns-3.44/src/vanet-security/model/vehicle-app.cc}

\lstinputlisting[
  style=cppstyle,
  firstline=410,
  lastline=518,
  caption={Stage 4 code: PBC signature generation.},
  label={lst:stage4-sign}
]{simulation_workspace/ns-allinone-3.44/ns-3.44/src/vanet-security/model/pbc-crypto.cc}

\lstinputlisting[
  style=cppstyle,
  firstline=520,
  lastline=612,
  caption={Stage 4 code: vehicle sends PBC or ECC V2V packet.},
  label={lst:stage4-send}
]{simulation_workspace/ns-allinone-3.44/ns-3.44/src/vanet-security/model/vehicle-app.cc}

\section{Stage 5: Relay Aggregation and Batch Verification}

\subsection{Goal}

The BS relay verifies a batch of PBC signatures before forwarding V2V packets.
This is the verification authority in LTE/NR relay mode.

\subsection{Formulas}

Aggregate signatures:

\[
\Psi=\sum_{i=1}^{n}\psi_i
\]

Aggregate pseudonym hashes:

\[
\sum\Pi=\sum_{i=1}^{n}H_1(\operatorname{PIDBytes}_i\parallel P_{pub})
\]

Aggregate public keys:

\[
\sum Q=\sum_{i=1}^{n}Q_i
\]

Accept batch if:

\[
e(\Psi,P)\stackrel{?}{=}
e(P_{pub},\sum\Pi)\cdot e(X,\sum Q)\cdot
\prod_{i=1}^{n} e(R_i,W_i)
\]

\subsection{Short Derivation}

From the signature formula:

\[
\psi_i=S_i+x_iX+w_iR_i
\]

Summing over all \(i\):

\[
\Psi=\sum_iS_i+\sum_ix_iX+\sum_iw_iR_i
\]

Using bilinearity and the definitions \(S_i=s\Pi_i\), \(P_{pub}=sP\),
\(Q_i=x_iP\), and \(W_i=w_iP\):

\[
e(\Psi,P)=e(P_{pub},\sum_i\Pi_i)\cdot e(X,\sum_iQ_i)\cdot\prod_i e(R_i,W_i)
\]

\subsection{Implementation}

\begin{itemize}[leftmargin=*]
  \item \codepath{BsRelayApp::HandleRead()} validates PBC fields and builds batches.
  \item \codepath{BsRelayApp::FlushBatch()} aggregates and verifies a batch.
  \item \codepath{PbcCrypto::AggregatePsi()} computes \(\Psi\).
  \item \codepath{PbcCrypto::VerifyAggregate()} implements the pairing equation.
\end{itemize}

\lstinputlisting[
  style=cppstyle,
  firstline=138,
  lastline=231,
  caption={Stage 5 code: relay validates packets and builds batches.},
  label={lst:stage5-batch}
]{simulation_workspace/ns-allinone-3.44/ns-3.44/src/vanet-security/model/bs-relay-app.cc}

\lstinputlisting[
  style=cppstyle,
  firstline=257,
  lastline=320,
  caption={Stage 5 code: relay flushes and verifies a batch.},
  label={lst:stage5-flush}
]{simulation_workspace/ns-allinone-3.44/ns-3.44/src/vanet-security/model/bs-relay-app.cc}

\lstinputlisting[
  style=cppstyle,
  firstline=530,
  lastline=705,
  caption={Stage 5 code: aggregate signature verification equation.},
  label={lst:stage5-verify}
]{simulation_workspace/ns-allinone-3.44/ns-3.44/src/vanet-security/model/pbc-crypto.cc}

\section{Evidence and Demo Checks}

For a clean LTE PBC smoke run, the logs and CSV should show:

\begin{lstlisting}[style=cmdstyle,caption={Useful professor-demo checks.}]
grep -c "TA PBC setup complete" /tmp/lte_pbc.log
grep -c "KGC PBC setup complete" /tmp/lte_pbc.log
grep -c "stored PID2" /tmp/lte_pbc.log
grep -c "stored partial key" /tmp/lte_pbc.log
grep -c "generated Qi" /tmp/lte_pbc.log
grep -c "missing active PBC signing material" /tmp/lte_pbc.log
awk -F, 'NR==2{print "verification_failures=" $4, "avg_kgc_rtt_ms=" $10, "avg_ta_rtt_ms=" $11}' results/lte_pbc.csv
\end{lstlisting}

Expected result for the successful smoke test:

\begin{itemize}[leftmargin=*]
  \item PBC setup messages appear at least once.
  \item Every registered vehicle stores PID2 and partial key.
  \item Every registered vehicle generates \(Q_i\).
  \item Missing active PBC signing material count is zero after registration completes.
  \item \codepath{verification_failures=0} in the CSV for a clean run.
\end{itemize}

\section{Important Limitations}

\begin{itemize}[leftmargin=*]
  \item PBC is implemented for the LTE/NR BS-relay path.
  \item Direct Wi-Fi/SUMO V2V remains on the ECC signing and verification path.
  \item ECC uses OpenSSL P-256. PBC uses the installed PBC library's Type A pairing parameters, not NIST P-256.
  \item Vehicles trust relay-forwarded PBC packets because the relay has already performed aggregate verification.
  \item Credential transport is simulation-level and is not protected by TLS.
\end{itemize}

\section{Compile Command}

Compile from the repository root so all listing paths resolve:

\begin{lstlisting}[style=cmdstyle,caption={Build this report locally.}]
pdflatex VANET_Security_Stages_Report.tex
pdflatex VANET_Security_Stages_Report.tex
\end{lstlisting}

\end{document}
